% Advanced Microeconomics II: Practice Examination on Game Theory and Information Economics
% Questions and answer key combined in a single document.
% Recommended compilation:
%   latexmk -pdf Micro2Q.tex

\documentclass[12pt,a4paper]{article}

% Packages
\usepackage[utf8]{inputenc}
\usepackage[margin=1in]{geometry}
\usepackage{amsmath}
\usepackage{amssymb}
\usepackage{amsfonts}
\usepackage{bm}
\usepackage{mathtools}
\usepackage{booktabs}
\usepackage{tabularx}
\usepackage{enumitem}
\usepackage{hyperref}

% Enumerate style for sub-questions
\setlist[enumerate,1]{label=\textbf{(\alph*)}, leftmargin=*}

\title{\textbf{Advanced Microeconomics II:}\\[0.25em]
       \Large Practice Examination on Game Theory and Information Economics}
\author{ }
\date{ }

\begin{document}

\maketitle

\section*{Examination Overview}

\begin{itemize}
  \item \textbf{Subject}: Game Theory and Information Economics
  \item \textbf{Time Allotted}: 3 Hours
  \item \textbf{Materials Allowed}: Non-programmable calculator.
  \item \textbf{Instructions}: Answer all four questions. Show all mathematical derivations rigorously.
        All equilibrium concepts must be clearly defined and justified using the standard axioms of
        non-cooperative game theory.
\end{itemize}

\bigskip
\noindent\hrulefill
\bigskip

%%%%%%%%%%%%%%%%%%%%%%%%%%%%%%%%%%%%%%%%%%%%%%%%%%%%%%%%%%%%%%%%%%%%%%%%%%%%%%%%
% 1. Simultaneous Games
%%%%%%%%%%%%%%%%%%%%%%%%%%%%%%%%%%%%%%%%%%%%%%%%%%%%%%%%%%%%%%%%%%%%%%%%%%%%%%%%

\section{Simultaneous Games: Dominance, Nash Equilibrium, and Security}

The strategic stability of simultaneous interactions is predicated on the internal consistency
of players' beliefs and the structural bounds of the strategic environment. In non-cooperative
settings, the \emph{Iterated Elimination of Strictly Dominated Strategies} (IESDS) provides a
mechanism for simplifying the strategy space by removing actions that are irrational under any
consistent belief system. Complementarily, the Maxmin criterion (security strategy) identifies the
foundational bounds of rational behavior, establishing the minimum utility a player can guarantee
themselves independently of the opponent's specific play---a concept of paramount importance in
environments where common knowledge of rationality may be absent.

\subsection*{Question 1: Strategic Reduction and Security Payoffs}

Consider the following \(4 \times 4\) payoff matrix for Player~A (rows) and Player~B (columns):

\begin{table}[h!]
  \centering
  \renewcommand{\arraystretch}{1.1}
  \begin{tabular}{lcccc}
    \toprule
    $A \setminus B$ & $b_1$ & $b_2$ & $b_3$ & $b_4$ \\
    \midrule
    $a_1$ & $(5,4)$ & $(4,5)$ & $(8,3)$ & $(3,2)$ \\
    $a_2$ & $(3,3)$ & $(1,2)$ & $(2,6)$ & $(1,1)$ \\
    $a_3$ & $(4,6)$ & $(2,4)$ & $(6,5)$ & $(2,3)$ \\
    $a_4$ & $(1,1)$ & $(0,2)$ & $(1,0)$ & $(2,1)$ \\
    \bottomrule
  \end{tabular}
\end{table}

\noindent\textbf{Tasks.}
\begin{enumerate}
  \item Perform the Iterated Elimination of Strictly Dominated Strategies (IESDS).
        Identify all strictly dominated strategies at each stage and state the final reduced game.
  \item Determine the Pure-Strategy Nash Equilibrium (PSNE) of the reduced game. Using the
        definition of Pareto efficiency, evaluate whether this equilibrium is efficient relative
        to the entire \(4 \times 4\) payoff matrix.
  \item Using the maxminimization procedure for mixed strategies (referencing the intersection
        of expected-utility planes as described in the course modules), calculate the security
        strategy and the security payoff \(v_i\) for both players in the original game.
\end{enumerate}

\bigskip
\noindent\textbf{Transition.}
While simultaneous games assume static timing, many industrial and strategic interactions involve
a clear leader--follower structure, necessitating the study of sequential commitment and the
refinement of Nash equilibria.

\bigskip
\noindent\hrulefill
\bigskip

%%%%%%%%%%%%%%%%%%%%%%%%%%%%%%%%%%%%%%%%%%%%%%%%%%%%%%%%%%%%%%%%%%%%%%%%%%%%%%%%
% 2. Sequential Games
%%%%%%%%%%%%%%%%%%%%%%%%%%%%%%%%%%%%%%%%%%%%%%%%%%%%%%%%%%%%%%%%%%%%%%%%%%%%%%%%

\section{Sequential Games: Subgame Perfection and Market Leadership}

The strategic advantage of commitment allows a first mover to fundamentally alter the follower's
incentives, thereby selecting a preferred equilibrium path. To ensure the credibility of such
commitments, we employ Subgame Perfect Nash Equilibrium (SPNE), which utilizes backward induction
to filter out non-credible threats---strategies that, while part of a Nash equilibrium in normal
form, would be suboptimal if the relevant subgame were actually reached.

\subsection*{Question 2: Asymmetric Stackelberg Duopoly and Niche Selection}

Consider a market with inverse demand
\[
  P(Q) = 150 - Q, \qquad Q = q_1 + q_2.
\]
The firms face asymmetric marginal costs \(c_1 = 10\) and \(c_2 = 25\). There are no fixed costs.

\begin{enumerate}
  \item Compute the Cournot--Nash equilibrium (simultaneous choice) quantities \((q_1^C, q_2^C)\)
        and the resulting profits \((\pi_1^C, \pi_2^C)\) for both firms.
  \item Derive the Stackelberg equilibrium where Firm~1 acts as the leader and Firm~2 as the
        follower. Calculate equilibrium quantities and profits. Quantify the ``first-mover
        advantage'' for Firm~1 relative to the Cournot benchmark.
  \item Introduce a \emph{niche choice} stage. Firm~1 and Firm~2 must sequentially choose between
        the \emph{Premium Niche} or the \emph{Mass Market} (Firm~1 chooses first). If both choose
        the same segment, they compete as Stackelberg duopolists as solved in part~(b). If they
        choose different segments, they act as monopolists in their respective niches, each earning
        a fixed niche profit \(V = 2000\). Entry into either segment requires a sunk cost
        \(E = 100\). Characterize the unique SPNE path and payoffs.
\end{enumerate}

\bigskip
\noindent\textbf{Transition.}
The structural stability of market entry often transitions into the complexities of bilateral
negotiation, where the timing of offers and the specificity of assets create new strategic
frictions.

\bigskip
\noindent\hrulefill
\bigskip

%%%%%%%%%%%%%%%%%%%%%%%%%%%%%%%%%%%%%%%%%%%%%%%%%%%%%%%%%%%%%%%%%%%%%%%%%%%%%%%%
% 3. Bargaining and Holdup
%%%%%%%%%%%%%%%%%%%%%%%%%%%%%%%%%%%%%%%%%%%%%%%%%%%%%%%%%%%%%%%%%%%%%%%%%%%%%%%%

\section{Bargaining and the Holdup Problem}

Incomplete contracting frequently results in the \emph{holdup problem}, wherein the anticipation of
future bargaining leads to ex-ante under-investment in relationship-specific assets. This
inefficiency arises because the investor, having sunk their costs, is vulnerable to a ``squeeze''
during the negotiation of the resulting surplus, failing to capture the full marginal return of
their investment.

\subsection*{Question 3: Alternating Offers and Specific Investment}

Consider a three-period alternating-offers bargaining game (\(T = 3\)) over a surplus \(S = 1\).
The discount factor is \(\delta \in (0,1)\). Player~1 makes offers in periods \(t = 1\) and
\(t = 3\); Player~2 makes the offer in period \(t = 2\).

\begin{enumerate}
  \item Solve for the unique SPNE shares of the surplus \((\phi_1, \phi_2)\) at \(t = 1\) using
        backward induction.
  \item Precede the bargaining with an investment stage. Player~2 can choose ``High'' investment
        (cost \(c_H = 0.2\), surplus \(V_H = 1.5\)) or ``Low'' investment
        (cost \(c_L = 0.05\), surplus \(V_L = 1.0\)). Given the bargaining shares derived in
        part~(a), determine Player~2's optimal investment choice.
  \item Evaluate whether a non-credible-threat Nash equilibrium exists that induces High
        investment. Explain, with reference to subgame perfection, why this threat fails to
        resolve the ex-ante inefficiency.
\end{enumerate}

\bigskip
\noindent\textbf{Transition.}
Strategic distortions are not limited to offer timing; they also emerge from the fundamental
inability of market participants to observe the hidden traits of the goods being traded.

\bigskip
\noindent\hrulefill
\bigskip

%%%%%%%%%%%%%%%%%%%%%%%%%%%%%%%%%%%%%%%%%%%%%%%%%%%%%%%%%%%%%%%%%%%%%%%%%%%%%%%%
% 4. Information Economics
%%%%%%%%%%%%%%%%%%%%%%%%%%%%%%%%%%%%%%%%%%%%%%%%%%%%%%%%%%%%%%%%%%%%%%%%%%%%%%%%

\section{Information Economics: Adverse Selection and Screening}

The ``lemons'' problem demonstrates how asymmetric information regarding quality (hidden traits)
can lead to adverse selection, where the market price reflects only the average quality of the
goods offered. If high-quality sellers cannot credibly signal their type, they may exit the market
entirely, resulting in market thinning or total collapse.

\subsection*{Question 4: Market Thinning and Thresholds}

Assume quality \(\theta\) is distributed uniformly on \([0,2]\). Seller valuation is
\[
  V_S(\theta) = \theta, \qquad
  V_B(\theta) = \theta + a, \quad a > 0.
\]

\begin{enumerate}
  \item If the buyer makes a take-it-or-leave-it price offer \(p\), solve for the buyer's optimal
        price \(p^*\). Determine the specific value of \(a\) for which the equilibrium probability
        of trade is exactly \(1/2\).
  \item Consider a competitive screening environment with multiple buyers. Derive the equilibrium
        price \(p_e\) and determine the condition on \(a\) required for all seller types
        \(\theta \in [0,2]\) to trade.
  \item In a discrete version, cars are either High Quality
        \((G, V_S = 16, V_B = 20)\) or Bad Quality \((B, V_S = 8, V_B = 10)\). Let \(\beta\) be
        the proportion of bad cars. Determine the threshold \(\beta^*\) at which High-Quality
        sellers exit the market under competitive equilibrium.
\end{enumerate}

\bigskip
\noindent\hrulefill
\bigskip

%%%%%%%%%%%%%%%%%%%%%%%%%%%%%%%%%%%%%%%%%%%%%%%%%%%%%%%%%%%%%%%%%%%%%%%%%%%%%%%%
% ANSWER KEY
%%%%%%%%%%%%%%%%%%%%%%%%%%%%%%%%%%%%%%%%%%%%%%%%%%%%%%%%%%%%%%%%%%%%%%%%%%%%%%%%

\section*{Answer Key}

This section provides one possible set of fully worked solutions. Unless otherwise stated,
all expectations and probabilities are taken under the distributions specified in the questions.

\subsection*{1. Simultaneous Games}

\subsubsection*{(a) Iterated Elimination of Strictly Dominated Strategies (IESDS)}

\paragraph{Round 1 (Row player A).}
Compare row \(a_2\) to \(a_1\):
\[
  (3,1,2,1) \quad \text{vs.} \quad (5,4,8,3).
\]
For Player~A's payoffs, we have \(3 < 5\), \(1 < 4\), \(2 < 8\), and \(1 < 3\). Thus \(a_1\) strictly
dominates \(a_2\) for Player~A. Eliminate strategy \(a_2\).

\paragraph{Round 2 (Column player B).}
In the reduced matrix with rows \(\{a_1,a_3,a_4\}\) and all four columns, compare columns \(b_4\)
and \(b_2\) for Player~B:
\[
  b_4: (2,3,1), \qquad b_2: (5,4,2).
\]
Coordinate-wise, \(2<5\), \(3<4\), and \(1<2\). Hence \(b_2\) strictly dominates \(b_4\) for
Player~B. Eliminate \(b_4\).

\paragraph{Round 3 (Row player A again).}
In the further-reduced game with rows \(\{a_1,a_3,a_4\}\) and columns \(\{b_1,b_2,b_3\}\),
compare \(a_4\) to \(a_1\):
\[
  a_4: (1,0,1), \qquad a_1: (5,4,8).
\]
Again \(1 < 5\), \(0 < 4\), and \(1 < 8\), so \(a_1\) strictly dominates \(a_4\) for Player~A.
Eliminate \(a_4\).

\paragraph{Reduced game.}
The final reduced game after IESDS is
\[
  \{a_1,a_3\} \times \{b_1,b_2,b_3\}.
\]

\subsubsection*{(b) Pure-Strategy Nash Equilibrium and Efficiency}

In the \(2 \times 3\) reduced game, compute each player's best responses.

\paragraph{Best responses for Player A.}
For each column \(b_j\), compare A's payoffs:
\begin{align*}
  b_1 &: \; u_A(a_1,b_1) = 5,\quad u_A(a_3,b_1) = 4
        &&\Rightarrow \text{BR}_A(b_1) = a_1, \\
  b_2 &: \; u_A(a_1,b_2) = 4,\quad u_A(a_3,b_2) = 2
        &&\Rightarrow \text{BR}_A(b_2) = a_1, \\
  b_3 &: \; u_A(a_1,b_3) = 8,\quad u_A(a_3,b_3) = 6
        &&\Rightarrow \text{BR}_A(b_3) = a_1.
\end{align*}

\paragraph{Best responses for Player B.}
For each row, compare B's payoffs:
\begin{align*}
  a_1 &: \; (b_1,b_2,b_3) = (4,5,3)
        &&\Rightarrow \text{BR}_B(a_1) = b_2, \\
  a_3 &: \; (b_1,b_2,b_3) = (6,4,5)
        &&\Rightarrow \text{BR}_B(a_3) = b_1.
\end{align*}

The unique mutual best response is
\[
  (a_1,b_2), \qquad \text{with payoffs } (4,5).
\]
Hence the unique pure-strategy Nash Equilibrium of the reduced game is \((a_1,b_2)\).

\paragraph{Pareto efficiency.}
In the original \(4 \times 4\) matrix, consider the outcome \((a_3,b_1)\) which yields
payoffs \((4,6)\). Relative to the equilibrium \((a_1,b_2) = (4,5)\), we have
\[
  4 \ge 4 \quad \text{and} \quad 6 > 5,
\]
so \((a_3,b_1)\) Pareto dominates \((a_1,b_2)\). Thus the Nash equilibrium is not Pareto efficient
in the original game.

\subsubsection*{(c) Security Strategies and Security Payoffs}

Restrict attention to the reduced game. Let Player~A mix between \(\{a_1,a_3\}\) with
\(\sigma_A = (p,1-p)\). Let Player~B mix among \(\{b_1,b_2,b_3\}\) with
\(\sigma_B = (q,1-q-r,r)\), and we will see that in the security strategy \(r = 0\).

\paragraph{Player A.}
For a given \(p\), A's payoffs against each pure column are
\begin{align*}
  u_A(\sigma_A,b_1) &= 5p + 4(1-p) = p + 4, \\
  u_A(\sigma_A,b_2) &= 4p + 2(1-p) = 2p + 2, \\
  u_A(\sigma_A,b_3) &= 8p + 6(1-p) = 2p + 6.
\end{align*}
The lower envelope is \(\min\{p+4,\,2p+2,\,2p+6\}\). For \(p \in [0,1]\), the binding (smallest)
expression is \(2p + 2\), which is maximized at \(p = 1\). Thus A's security strategy is the pure
strategy \(a_1\), and her security payoff is
\[
  v_A = 4.
\]

\paragraph{Player B.}
Let B mix over \(b_1\) and \(b_2\) only with probabilities \((q,1-q,0)\). Then, for each row:
\begin{align*}
  u_B(a_1,\sigma_B) &= 4q + 5(1-q) = 5 - q, \\
  u_B(a_3,\sigma_B) &= 6q + 4(1-q) = 2q + 4.
\end{align*}
To maximize the minimum payoff, equalize these:
\[
  5 - q = 2q + 4
  \quad\Longrightarrow\quad
  3q = 1
  \quad\Longrightarrow\quad
  q = \tfrac{1}{3}.
\]
The associated security payoff is
\[
  v_B = 5 - q = 5 - \tfrac{1}{3} = \tfrac{14}{3} \approx 4.67.
\]
Thus B's security strategy is
\[
  \sigma_B = \left(\tfrac{1}{3}, \tfrac{2}{3}, 0, 0\right).
\]

\bigskip

\subsection*{2. Sequential Games}

\subsubsection*{(a) Cournot--Nash Equilibrium}

Firm \(i\)'s profit is
\[
  \pi_i(q_1,q_2) = \bigl(150 - q_1 - q_2 - c_i\bigr)\, q_i.
\]
Best-response functions are
\begin{align*}
  \text{BR}_1(q_2) &= \frac{150 - q_2 - c_1}{2}
                   = 70 - \tfrac{1}{2}q_2, \\
  \text{BR}_2(q_1) &= \frac{150 - q_1 - c_2}{2}
                   = 62.5 - \tfrac{1}{2}q_1.
\end{align*}

Solving for the Cournot equilibrium:
\[
  q_1 = 70 - \tfrac{1}{2}(62.5 - \tfrac{1}{2}q_1)
  \quad\Longrightarrow\quad
  0.75\,q_1 = 38.75
  \quad\Longrightarrow\quad
  q_1^C \approx 51.67,
\]
and
\[
  q_2^C = 62.5 - \tfrac{1}{2}q_1^C \approx 36.67.
\]
Price is
\[
  P^C = 150 - q_1^C - q_2^C \approx 61.67.
\]
Profits:
\begin{align*}
  \pi_1^C &\approx (61.67 - 10)\cdot 51.67 \approx 2669.7,\\
  \pi_2^C &\approx (61.67 - 25)\cdot 36.67 \approx 1344.7.
\end{align*}

\subsubsection*{(b) Stackelberg Equilibrium with Firm 1 as Leader}

In Stackelberg, Firm~1 chooses \(q_1\) anticipating that Firm~2 plays its Cournot best response:
\[
  q_2(q_1) = 62.5 - \tfrac{1}{2}q_1.
\]
Thus total quantity is
\[
  Q = q_1 + q_2(q_1) = 62.5 + \tfrac{1}{2}q_1,
\]
and price is
\[
  P = 150 - Q = 87.5 - \tfrac{1}{2}q_1.
\]
Firm~1's profit as a function of \(q_1\) is
\[
  \pi_1(q_1) = (P - c_1)q_1
             = (87.5 - \tfrac{1}{2}q_1 - 10)q_1
             = (77.5 - \tfrac{1}{2}q_1)q_1.
\]
The FOC gives
\[
  \frac{d\pi_1}{dq_1} = 77.5 - q_1 = 0
  \quad\Longrightarrow\quad
  q_1^S = 77.5.
\]
Then
\[
  q_2^S = 62.5 - \tfrac{1}{2}q_1^S = 62.5 - 38.75 = 23.75.
\]
Stackelberg profits:
\begin{align*}
  \pi_1^S &= (77.5 - \tfrac{1}{2}\cdot 77.5)\cdot 77.5
          = 38.75 \cdot 77.5
          \approx 3003.1, \\
  \pi_2^S &= (P^S - c_2)\,q_2^S
          \approx (150 - 101.25 - 25)\cdot 23.75
          \approx 564.1.
\end{align*}
The first-mover advantage for Firm~1 is
\[
  \pi_1^S - \pi_1^C \approx 3003.1 - 2669.7 = 333.4.
\]

\subsubsection*{(c) Niche Choice SPNE}

If both firms choose the same segment (Premium or Mass), they play the Stackelberg duopoly from
part~(b) and, net of entry costs,
\[
  \pi_1^{\text{same}} = \pi_1^S - E \approx 3003.1 - 100 = 2903.1,\qquad
  \pi_2^{\text{same}} = \pi_2^S - E \approx 564.1 - 100 = 464.1.
\]
If they choose different segments, each is a monopolist in their own niche and obtains
\[
  \pi_i^{\text{diff}} = V - E = 2000 - 100 = 1900.
\]

We analyze the niche-choice stage by backward induction. Suppose Firm~1 has chosen Mass:
Firm~2 can either also choose Mass (yielding \(464.1\)) or choose Premium (yielding \(1900\)).
Thus its best response is to pick the opposite niche. Analogously, if Firm~1 picks Premium,
Firm~2's best response is to choose Mass.

Anticipating this, Firm~1 understands that choosing either niche leads to the outcome where
the other firm chooses the opposite niche and both obtain \(1900\). Hence any initial niche
choice by Firm~1 is optimal (they are indifferent), and Firm~2's unique best response is to
choose the other segment.

Therefore, the unique SPNE path features the firms in different niches, and the final payoffs are
\[
  (\pi_1,\pi_2) = (1900,1900).
\]

\bigskip

\subsection*{3. Bargaining and the Holdup Problem}

\subsubsection*{(a) SPNE Shares in the \(T=3\) Bargaining Game}

Work backward from period \(t=3\).
\begin{itemize}
  \item \textbf{Period 3.} Player~1 makes the final offer. Any accepted offer must give Player~2
        at least \(0\). Player~1 can therefore extract the entire surplus, offering
        \((1,0)\). Thus the continuation value at \(t=3\), for Player~1, is \(1\).
  \item \textbf{Period 2.} Player~2 makes an offer \((x_2,1-x_2)\) knowing that rejection leads to
        period~3 where Player~1 obtains \(1\). Hence Player~1's continuation value from rejection
        at \(t=2\) is \(\delta \cdot 1\). To secure acceptance, Player~2 must give Player~1 at least
        \(\delta\), so the best she can do is to offer \((\delta,1-\delta)\). Thus the period-2
        continuation value for Player~2 is \(1-\delta\).
  \item \textbf{Period 1.} Player~1 now proposes \((x_1,1-x_1)\). If Player~2 rejects, the game
        moves to period~2 where Player~2's expected payoff is \((1-\delta)\), discounted once:
        \(\delta(1-\delta)\). To induce acceptance, Player~1 must give Player~2 at least this
        amount. Hence the profit-maximizing offer is
        \[
          (\phi_1,\phi_2) = \bigl(1-\delta+\delta^2,\; \delta-\delta^2\bigr).
        \]
\end{itemize}

\subsubsection*{(b) Optimal Investment Choice for Player 2}

If Player~2 chooses High investment, the surplus becomes \(V_H = 1.5\), so her payoff is
\[
  \pi_2(H) = 1.5(\delta - \delta^2) - 0.2.
\]
If she chooses Low investment, the surplus is \(V_L = 1.0\), and her payoff is
\[
  \pi_2(L) = 1.0(\delta - \delta^2) - 0.05.
\]
High investment is optimal if and only if
\[
  \pi_2(H) > \pi_2(L)
  \quad\Longleftrightarrow\quad
  1.5(\delta - \delta^2) - 0.2 > (\delta - \delta^2) - 0.05,
\]
which simplifies to
\[
  0.5(\delta - \delta^2) > 0.15
  \quad\Longleftrightarrow\quad
  \delta - \delta^2 > 0.3.
\]
Thus Player~2 chooses High investment whenever \(\delta - \delta^2 > 0.3\).

\subsubsection*{(c) Non-credible Threats and Ex-ante Inefficiency}

One can construct a Nash equilibrium in which Player~1 \emph{promises} an offer in the bargaining
subgame that leaves Player~2 with a sufficiently high share (say, \(x > \delta - \delta^2\)) to
induce High investment ex ante. However, this threat is not credible.

Once the investment cost is sunk and the game enters the bargaining subgame, Player~1's dominant
action is to play the SPNE strategy derived in part~(a), namely to offer
\[
  (1 - \delta + \delta^2,\; \delta - \delta^2),
\]
which maximizes his continuation payoff. Any strategy profile in which Player~1 deviates from this
offer in the subgame is not sequentially rational and hence cannot be part of a Subgame Perfect
Nash Equilibrium. Therefore, such promises cannot solve the ex-ante under-investment problem:
the holdup problem persists.

\bigskip

\subsection*{4. Information Economics: Adverse Selection and Screening}

\subsubsection*{(a) Take-it-or-leave-it Offer and Trade Probability}

For a given price \(p\), only sellers with \(\theta \le p\) are willing to trade, since
their valuation is \(V_S(\theta) = \theta\). Under a uniform distribution on \([0,2]\), the
probability of trade is
\[
  \Pr(\theta \le p) = \frac{p}{2}, \quad 0 \le p \le 2.
\]
Conditional on trade, the expected quality is
\[
  \mathbb{E}[\theta \mid \theta \le p] = \frac{p}{2}.
\]
Thus the buyer's expected valuation is
\[
  \mathbb{E}[V_B(\theta) \mid \theta \le p]
  = \mathbb{E}[\theta + a \mid \theta \le p]
  = \frac{p}{2} + a.
\]
The buyer's expected payoff is
\[
  U_B(p) = \Pr(\theta \le p)\,\bigl(\mathbb{E}[V_B \mid \theta \le p] - p\bigr)
         = \frac{p}{2} \left(\frac{p}{2} + a - p\right)
         = \frac{ap}{2} - \frac{p^2}{4}.
\]
Maximizing with respect to \(p\):
\[
  \frac{dU_B}{dp} = \frac{a}{2} - \frac{p}{2} = 0
  \quad\Longrightarrow\quad
  p^* = a \quad (\text{for } a \le 2).
\]
The associated probability of trade is
\[
  \Pr(\theta \le p^*) = \frac{p^*}{2} = \frac{a}{2}.
\]
For the equilibrium probability of trade to equal \(1/2\), we require
\[
  \frac{a}{2} = \frac{1}{2} \quad\Longrightarrow\quad a = 1.
\]

\subsubsection*{(b) Competitive Screening Equilibrium}

Under competition between many buyers, the price \(p_e\) must equal the expected buyer valuation
conditional on trade:
\[
  p_e = \mathbb{E}[V_B(\theta) \mid \theta \le p_e]
      = \frac{p_e}{2} + a.
\]
Solving,
\[
  p_e - \frac{p_e}{2} = a
  \quad\Longrightarrow\quad
  \frac{p_e}{2} = a
  \quad\Longrightarrow\quad
  p_e = 2a.
\]
For all seller types \(\theta \in [0,2]\) to be willing to trade, we need
\[
  p_e \ge 2
  \quad\Longrightarrow\quad
  2a \ge 2
  \quad\Longrightarrow\quad
  a \ge 1.
\]

\subsubsection*{(c) Discrete Threshold with Good and Bad Cars}

Let \(\beta\) be the proportion of bad cars (type \(B\)), and \(1-\beta\) be the proportion of good
cars (type \(G\)). Buyer valuations are
\[
  V_B(G) = 20, \qquad V_B(B) = 10.
\]
In competitive equilibrium, the price equals the expected buyer valuation:
\[
  p = \mathbb{E}[V_B]
    = \beta \cdot 10 + (1-\beta)\cdot 20
    = 20 - 10\beta.
\]
High-quality sellers will remain in the market only if the price at least covers their valuation
of \(16\):
\[
  p \ge 16
  \quad\Longrightarrow\quad
  20 - 10\beta \ge 16
  \quad\Longrightarrow\quad
  4 \ge 10\beta
  \quad\Longrightarrow\quad
  \beta \le 0.4.
\]
Thus the threshold is
\[
  \beta^* = 0.4.
\]

\end{document}

